% Generated draft: requires mathematical and source review.
\documentclass[11pt]{article}
\usepackage[margin=25mm]{geometry}
\usepackage{fontspec}
\setmainfont{Latin Modern Roman}
\usepackage{amsmath,amssymb,mathtools,unicode-math}
\setmathfont{Latin Modern Math}
\usepackage{enumitem,needspace,xurl,hyperref}
\setlength{\parindent}{0pt}
\setlength{\parskip}{5pt}
\setlength{\emergencystretch}{3em}
\newcommand{\N}{\mathbb{N}}
\newcommand{\Z}{\mathbb{Z}}
\newcommand{\Q}{\mathbb{Q}}
\newcommand{\R}{\mathbb{R}}
\newcommand{\C}{\mathbb{C}}
\newcommand{\F}{\mathbb{F}}
\newcommand{\A}{\mathbb{A}}
\newcommand{\PP}{\mathbb{P}}
\newcommand{\E}{\mathbb{E}}
\newcommand{\eps}{\varepsilon}
\newcommand{\dd}{\,\mathrm{d}}
\begin{document}
\section*{UnsolvedMath ID: 700404}
\textbf{Source problem code:} catalogue entry
\subsection{Definitions and mathematical statement}
Let \(G\) be a group and let \(p\) be a prime. A finite maximal subgroup of \(G\) is a subgroup \(M\) satisfying
\[
|M|<\infty,\qquad M<G,\qquad
\forall H\leq G\;\bigl(M\leq H\Longrightarrow H=M\text{ or }H=G\bigr).
\]
A \(p\)-group is a group whose elements all have finite order a power of \(p\). A maximal \(p\)-subgroup of \(G\) is a \(p\)-subgroup \(P\leq G\) such that every \(p\)-subgroup \(Q\leq G\) containing \(P\) equals \(P\). This differs from a maximal proper subgroup of \(G\) which happens to be a \(p\)-group. A periodic group is a group in which every element has finite order.

The source asks whether a ``(conjugacy, \(p\)-conjugacy) biprimitively finite'' group must be finite if it has a finite maximal subgroup ``(\(p\)-subgroup).'' For the finite maximal subgroup formulation, the requested implication has the form
\[
\forall G\;\left[
\left(G\text{ has the intended biprimitive finiteness property}\right)
\land
\left(\exists M<G:\ |M|<\infty\text{ and }M\text{ is maximal in }G\right)
\Longrightarrow |G|<\infty
\right].
\]
Here the phrase ``intended biprimitive finiteness property'' is a placeholder for the source's terminology, not a definition. The supplied record does not define biprimitive finiteness, conjugacy biprimitive finiteness, or \(p\)-conjugacy biprimitive finiteness. Consequently, this implication is a formulation schema requiring those definitions before it is a self-contained mathematical statement.

The parenthetical \(p\)-subgroup formulation also requires clarification: it may replace the subgroup hypothesis by the existence of a finite maximal \(p\)-subgroup, or require a finite maximal proper subgroup which is a \(p\)-group. The supplied record does not explicitly resolve this distinction or specify how its parenthetical property variants correspond to its subgroup variants. No assertion that every combination is intended is made here.

The source reports affirmative answers for conjugacy biprimitively finite \(p\)-groups and for \(2\)-conjugacy biprimitively finite groups, attributing these results to V. P. Shunkov. It also reports that the statement fails for arbitrary periodic groups and refers to Archive, Problem 3.9. These are preserved source qualifications; their precise hypotheses and correspondence with the variants above require checking the cited material.
\subsection{Short English statement}
Problem 6.62, proposed by V. P. Shunkov, asks whether the indicated forms of biprimitive finiteness force a group with a finite maximal subgroup, or the intended \(p\)-subgroup variant, to be finite. The record reports two affirmative special cases and failure for arbitrary periodic groups. Definitions of the specialized finiteness properties and the interpretation of the parentheses remain unresolved.
\subsection{Sources}
\begin{itemize}
\item Supplied local UnsolvedMath registry record: permanent ID \(700404\), exact source problem code \texttt{catalogue entry}, title \emph{Kourovka Notebook Problem 6.62}. The registry labels the problem \texttt{open}; this is a supplied status, not an independently verified assessment.
\item \emph{The Kourovka Notebook (Unsolved Problems in Group Theory)}, Problem 6.62, proposed by V. P. Shunkov; sixth issue (1978), as recorded in the supplied twenty-first-issue edition (2026), \texttt{arXiv:1401.0300v48}. The supplied collection description names E. I. Khukhro and V. D. Mazurov as editors.
\item V. P. Shunkov, \emph{Algebra and Logic} \textbf{9}, no. 4 (1970), 291--297; \textbf{11}, no. 4 (1972), 260--272; \textbf{12}, no. 5 (1973), 347--353. The source cites these collectively for the reported affirmative special cases; article titles and a correspondence between individual articles and cases are not supplied.
\item \emph{The Kourovka Notebook}, Archive, Problem 3.9: supplied cross-reference for the reported failure for arbitrary periodic groups. Its contents are not included in the supplied record.
\end{itemize}
\end{document}
